% =====================================================================
% VERSION ENSAE, REDUCTION SFE : chapitre 2, revue de littérature et positionnement.
% Phase 3 : copie a texte constant de chapitres/03_etat_art_positionnement.tex.
% Phase 4 (1er octobre 2026) : reecrit et condense ; toutes les citations, le tableau
% comparatif et les nombres sont conserves. Le texte d'origine reste dans chapitres/
% (version Institut) et dans l'historique Git.
% Lu par main_ensae.tex seul : ne jamais le faire lire par main_institut.tex,
% qui garde chapitres/.
% =====================================================================
\chapter{Revue de littérature et positionnement}
\label{ch:positionnement}

% SOURCES-SCRIPTS: 08g 32 01 23

La modélisation actuarielle du risque cyber s'est construite en trois vagues. La première traite
l'assurabilité du risque et la rareté des données \citep{BienerElingWirfs2015}, et décrit les
marges par la théorie des valeurs extrêmes. La deuxième modélise la contagion entre entités d'un
portefeuille par des modèles de réseau ou épidémiologiques, et la troisième décrit la dynamique des
arrivées par des processus de Hawkes.

Nous reprenons à ces travaux les marges empiriques et le cadre du risque extrême. Au risque de
crédit, nous empruntons le processus de branchement multitype pour la propagation entre piliers
(chapitre~\ref{chap:cascade}), et les modèles multi-états de migration de rating pour l'état de
conformité (chapitre~\ref{chap:multietats}). 

\section{Dépendance dirigée et dépendance à l'ordre}
% SOURCES-SCRIPTS: 08g 32 01 23

La matrice $W$ de contagion entre piliers (posée à dire d'expert, puis bornée au chapitre~\ref{chap:partielle}) donne par $W_{jk}$ la propension d'une défaillance de $j$ à entraîner celle de $k$. La ligne désigne l'émetteur, la colonne le récepteur. Elle est
dirigée, au sens où $W_{jk}\neq W_{kj}$, et autorise des cycles comme P1$\to$P2$\to\cdots\to$P1. La
normalisation de Leontief permet de stabiliser ces cycles ($\rho(W)\le g<1$, forme réduite $(I-W)^{-1}$).

Avec le score de criticité du modèle, de $0$ à $10$ et construit sur $W$ (chapitre~\ref{chap:cascade}), $22$ des $26$ sous-ensembles d'au moins deux piliers changent de score selon l'ordre. L'écart atteint $6$ points, entre $8$ pour la chaîne P1$\to$P2 et $2$ pour la chaîne
P2$\to$P1. Un score statique invariant par permutation ne peut pas
reproduire cette propriété. Nous avons aussi testé, sous trois formulations, une non-transitivité de la dominance entre piliers ; elle ne se vérifie pas et n'est pas revendiquée (chapitre~\ref{chap:cascade}).

\section{Comparaison aux approches existantes}
% SOURCES-SCRIPTS: 32 23 01

\begin{center}\footnotesize
\renewcommand{\arraystretch}{1.35}
\begin{tabular}{@{}p{3.1cm} c c c c p{3.1cm} p{2.9cm}@{}}
\toprule
\textbf{Approche} & \textbf{Dirigée} & \textbf{Ordre} & \textbf{Cycles} & \textbf{Propagation} &
\textbf{Donnée requise} & \textbf{Objet modélisé} \\
\midrule
Matrices de risque / AMDEC & \non & \non & \non & \non & dires d'expert & risques listés (probabilité $\times$ gravité) \\
Copule (Gumbel, gaussienne) & \non & \non & \non & \non & co-mouvement observé & briques de perte \\
Réseau / épidémiologie (SIR) & \oui & \non & \oui & \oui & structure réseau $+$ taux d'infection & entités d'un portefeuille \\
Hawkes (multivarié) & \oui & \non & \oui & \oui & horodatage fin des arrivées & types d'attaque, dans le temps \\
Réseau bayésien & \oui & \non & \non & \oui & tables de probabilités conditionnelles & variables d'un graphe acyclique \\
\cascell\textbf{Cascade dirigée (ici)} & \cascell\oui & \cascell\oui & \cascell\oui & \cascell\oui & \cascell jugement d'expert, puis bornes & \cascell\textbf{domaines de contrôle (piliers) d'une entité} \\
\bottomrule
\end{tabular}
\end{center}

\vspace{2pt}
\noindent Les matrices de risque et l'analyse des modes de défaillance, de leurs effets et de leur
criticité (AMDEC), standard des praticiens, notent des risques sans propagation ni direction.
L'AMDEC les classe par un indice de priorité (RPN $=$ gravité $\times$ occurrence $\times$ détection). Les copules
\citep{BohmeKataria2006, HerathHerath2011} décrivent une dépendance statique entre risques, sans ordre d'atteinte. Dans la comparaison du chapitre~\ref{chap:resultats}, leur interaction moyenne entre piliers est nulle par construction ($+0$), contre $+560$~M\euro{} (périmètre financier étudié) pour la cascade. Leur ordre de priorité diffère de celui de la cascade sur P2 et P3 (chapitre~\ref{chap:resultats}).

Les modèles de réseau et épidémiologiques \citep{HillairetLopez2021,
HillairetLopezDOultremontSpoorenberg2022, XuHua2017, Baldwin2017} modélisent la contagion entre
entités d'un portefeuille et exigent la structure du réseau et des taux d'infection, indisponibles
ici. Nous reprenons cependant à Hillairet et Lopez l'idée d'états successifs, transposée à l'état de conformité de chaque pilier. Les processus de Hawkes \citep{BessyRolandBoumezouedHillairet2021, BoumezouedCherkaouiHillairet2023}
décrivent l'auto-excitation entre types d'attaque. Sur le nombre annuel d'incidents, les données ne séparent pas l'auto-excitation d'une co-occurrence exogène le même jour (chapitre~\ref{ch:robustesse}). Nous retenons la seconde lecture, plus parcimonieuse ; elle se distingue du modèle à deux phases de $2023$, qui conserve l'auto-excitation tout en ajoutant une excitation externe. Les réseaux bayésiens intègrent la direction, mais supposent un graphe acyclique
et des probabilités conditionnelles qui se heurtent au même problème d'identification que $W$
(chapitre~\ref{ch:identifiabilite}).

\section{Travaux par thème}

\subsection*{Dépendance et contagion}
% SOURCES-SCRIPTS: 08g
La corrélation entre risques cyber a d'abord été traitée par
copules \citep{BohmeKataria2006, HerathHerath2011}. Sur la chronologie PRC, l'indice de Gini de la charge journalière vaut $0{,}69$ ($2\,922$ jours, $4\,614$ événements) : les sinistres se concentrent sur peu de jours, ce qui est compatible avec une accumulation par cause commune.
\citet{HillairetLopez2021} couplent des processus de comptage à un modèle épidémiologique
compartimental de type SIR (susceptibles, infectés, rétablis), et \citet{HillairetLopezDOultremontSpoorenberg2022} y ajoutent une
structure de réseau. \citet{XuHua2017} et \citet{Baldwin2017} modélisent la contagion entre cibles.

\subsection*{Dynamique temporelle}
\citet{BessyRolandBoumezouedHillairet2021} ajustent des processus de
Hawkes aux données PRC, l'une de nos sources, puis \citet{BoumezouedCherkaouiHillairet2023}
ajoutent un terme exogène. \citet{HillairetReveillacRosenbaum2021} en développent la théorie pour le
pricing de dérivés, et \citet{Peng2016} les appliquent à des taux d'attaque horodatés.

\subsection*{Fréquence, sévérité et queue}
\citet{Edwards2016} ajustent une binomiale négative au
nombre journalier de brèches, et \citet{LuZhangZhu2024} séparent fréquence et nombre de brèches par
événement, une structure compatible avec la loi de fréquence retenue ici (binomiale négative, chapitre~\ref{chap:socle}). \citet{FarkasLopezThomas2021}
tarifent par arbres de régression GPD. La queue lourde des pertes cyber est documentée
\citep{MaillartSornette2010, wheatley2016}, et \citet{elingloperfido2017} comparent plusieurs
familles de lois sur les brèches de données. \citet{elingwirfs2019} étudient les pertes cyber de la
base SAS OpRisk Global par les valeurs extrêmes, et leur étude est le précédent le plus proche pour la sévérité.

\subsection*{Travaux sur les mêmes sources}
\citet{Ranguin2024} modélise le risque cyber des
entreprises françaises pour la Caisse Centrale de Réassurance, à partir de PRC, VCDB, Hackmageddon et
de l'étude LUCY. La sévérité y est construite par arbres de régression GPD sur le volume de données
compromises. La sinistralité d'un portefeuille fictif est simulée, sinistres courants compris. Les lois
sont ensuite reparamétrées dans un cadre bayésien sur le marché français. Cette étude ne produit pas de
mesure de capital à un quantile réglementaire. Notre objet est le besoin de capital d'une
entité financière selon son état de conformité. Nous calibrons la sévérité sur des pertes en montants et
modélisons la dépendance entre piliers.

Sur PRC, \citet{KherLopezRapior2023} estiment la gravité d'un incident à partir de sa description
textuelle par réseaux de neurones, et \citet{Dountio2026} prolonge cette démarche. Ces travaux
classent la sévérité d'un incident en volume de dossiers compromis.

\subsection*{Cadre structurel et identification partielle}
Le facteur latent commun vient du modèle à facteur de \citet{Vasicek2002}, et la normalisation qui stabilise une contagion avec
cycles provient de l'analyse entrées-sorties de \citet{Leontief1936}.

Le traitement du paramètre non identifiable relève de la littérature économétrique des bornes. On y
caractérise l'ensemble des valeurs compatibles avec les données, et l'on en déduit un intervalle pour
la quantité d'intérêt. Initiée par les bornes de sélection de \citet{Manski1990}, exposée dans
\citet{Manski1995} et systématisée par \citet{Tamer2010}, cette approche est courante en évaluation
des politiques publiques. Elle est presque absente de l'actuariat du cyber, où l'on pose d'ordinaire
un paramètre non observable à dire d'expert. Nous publions le capital en bornes sur l'ensemble admissible
(chapitre~\ref{chap:partielle}). Une échelle de repli sur l'écart $u_{ij}$ à la dépendance moyenne y gradue les hypothèses de
structure.

\subsection*{Cascade de risques climatiques}
Un préprint de 2026 propose un réseau de cascade pour le
risque climatique de catastrophe \citep{ccrn2026}, sans péril, données ni auteurs communs avec ce
travail. Il s'agit d'un graphe dirigé de propagation entre n\oe uds, avec des indices de temps
séparés. Ce préprint est synthétique, et ses auteurs mettent en garde contre toute lecture de ses
valeurs comme calibrées. Nous n'en reprenons aucun chiffre ; nous empruntons l'architecture : séparation d'échelles, protocole de comparaison, expérience d'ablation et forme
fonctionnelle.

\section{Positionnement}

Nous reprenons la queue de sévérité par les valeurs extrêmes, la structure de paquet de la fréquence
et l'idée d'états successifs de Hillairet et Lopez, transposée à la conformité. Deux éléments sont nouveaux : une
dépendance dirigée et à l'ordre entre les piliers DORA, traduite en capital, et une frontière
d'identifiabilité explicite, traduite en recommandation sur les données à déclarer. Aucune des références que nous avons lues ne modélise une dépendance dirigée entre piliers d'une entité ni ne la traduit en besoin de capital (ORSA). Les lois de sévérité et de fréquence sont reprises de la littérature et calibrées.

\nocite{BohmeKataria2006, HerathHerath2011, HillairetLopez2021,
  HillairetLopezDOultremontSpoorenberg2022, XuHua2017, Baldwin2017,
  BessyRolandBoumezouedHillairet2021, BoumezouedCherkaouiHillairet2023,
  HillairetReveillacRosenbaum2021, Peng2016, Edwards2016, LuZhangZhu2024,
  FarkasLopezThomas2021, Vasicek2002, Leontief1936}
